\documentclass[11pt]{article}

% arXiv preprint: non-anonymous, with page numbers.
\usepackage[preprint]{acl}

\usepackage{times}
\usepackage{latexsym}
\usepackage[T1]{fontenc}
\usepackage[utf8]{inputenc}
\usepackage{amsmath}
\usepackage{microtype}
\usepackage{inconsolata}
\usepackage{graphicx}
\usepackage{booktabs}
\usepackage{multirow}

\graphicspath{{../results/figures/}}

% fig6 (the banded winner map) is a tall single-column float. At LaTeX's
% defaults it does not clear \floatpagefraction and is deferred past the
% bibliography. Letting a float page be mostly float puts it back in the body,
% beside the text that refers to it.
\renewcommand{\topfraction}{0.92}
\renewcommand{\textfraction}{0.07}
\renewcommand{\floatpagefraction}{0.72}
\setcounter{topnumber}{2}
\setcounter{totalnumber}{3}
\setlength{\textfloatsep}{10pt plus 2pt minus 2pt}
\setlength{\intextsep}{10pt plus 2pt minus 2pt}
\setlength{\abovecaptionskip}{4pt}

\newcommand{\us}{\ensuremath{\mu}s}
\newcommand{\code}[1]{\texttt{\small #1}}

\title{Attention Kernel Portability Across Six GPUs and Two Inference Regimes:\\
       A Measurement Study}

\author{Srivatsan Sarvesan \\
  University of California, San Diego \\
  \texttt{ssarvesan@ucsd.edu} \\}

\begin{document}
\maketitle

\begin{abstract}
Published attention benchmarks report a winner measured on one GPU. Does that
winner survive a change of device? We time 10 prefill and 6 decode
implementations over 403 configurations on six NVIDIA GPUs spanning sm80, sm86,
sm89, sm90 and sm120, under one software distribution: torch 2.9.0+cu128, CUDA
12.8, triton 3.5.0, flash-attn 2.8.3, flashinfer 0.6.18. Most configurations have no winner worth reporting: the fastest backend is
separated from the runner-up, by both a 10\% margin and a bootstrap interval
excluding 1, on 29\% of 447 forward-prefill and 34\% of 1203 decode
configurations. Among those that are separated, the fastest backend changes
between devices on 27 of 64 forward-prefill comparisons and 29 of 289 decode
comparisons within Ampere, Ada and Hopper. Changing winner is not the same as
paying for it: transferring the source device's choice costs a median
1.002$\times$ in prefill and 1.000$\times$ in decode, with 95th percentiles of
1.61$\times$ and 1.32$\times$, and 14\% of source choices are unavailable on
the target at all. Availability is where the comparison breaks down, at 63 of 85 decode comparisons against the RTX 5090, and availability rather than
architecture accounts for 62 of those 64: FlashInfer, which wins 131 of 244
decode configurations on H100, produced no timing at all on the RTX 5090,
raising a compute-capability \code{RuntimeError} on 950 of 960 attempted rows. A
cuobjdump audit of the wheels shows the same asymmetry statically: flash-attn
2.8.3 ships no sm\_86 or sm\_89 cubins and no PTX, so on those parts it runs
sm\_80 code. The unit of performance is not (kernel, GPU) but (kernel, GPU,
regime, software-distribution snapshot).
\end{abstract}

\section{Introduction}

A practitioner choosing an attention kernel reads a benchmark run on one GPU.
Does its winner survive a change of device?

The question is usually assumed rather than measured. Serving systems fix an
attention backend by preference order and carry that order onto every device
they deploy on, and kernel papers report a winner on the one part their authors
had. If the ordering belongs to the kernel, that practice is sound. If it
belongs to the pairing of kernel and device, a preference order tuned on an
A100 is a guess everywhere else.

We time 10 prefill and 6 decode implementations over 403 configurations on the
six GPUs of Table~\ref{tab:devices}; 73{,}230 measured rows collapse to
11{,}858 usable (GPU, configuration, implementation) medians.\footnote{Harness, analysis and figure code: \url{https://github.com/Srivatsan6923/attention-kernel-portability}}
Every claim is a within-device ordering compared across devices, never a claim
that one device is faster than another, so host effects such as clock drift and
contention are common to the implementations being compared.

Four results follow. First, most configurations have no decisive winner: on
68.1\% of them the fastest implementation is not separated from the runner-up
by both a 10\% margin and a bootstrap interval excluding 1, and a winner
reported there did not meet our separation criterion. Second, among configurations that do have a decisive
winner, decode orderings largely transfer across Ampere, Ada and Hopper while
prefill orderings do not. Third, decode agreement collapses against the
Blackwell part, and availability is where the comparison breaks down: in most of
those flips the winner on the other device is an implementation that produced no
timing at all on sm120.
Fourth, an audit of the shipped binaries shows the same asymmetry statically,
before anything is timed.

Two contributions follow from that. The first is a protocol that records which
kernel actually ran rather than which was requested, and that reports a winner
only where it is separated from the runner-up. The second is the finding itself:
the unit of performance is not (kernel, GPU) but (kernel, GPU, regime,
software-distribution snapshot).

\section{Background and related work}

Each attention paper names one device. FlashAttention
\citep{dao2022flashattention} keeps the $N \times N$ scores out of HBM, on A100;
FlashAttention-2 \citep{dao2023flashattention2} reaches 50--73\% of peak FLOPs/s,
again on A100, deferring Hopper; FlashAttention-3
\citep{shah2024flashattention3} is that deferred work, on H100 only, at 840 BF16
TFLOPs/s. Flash-Decoding \citep{dao2023flashdecoding} adds split-KV for
$q_{\text{len}} = 1$ on A100, and FlashInfer \citep{ye2025flashinfer} generalises
it into a serving engine, on A100 and H100. The generation-spanning reading is
the reader's addition, not the authors'.

Evaluation inherits the habit. The PyTorch 2 paper \citep{ansel2024pytorch2}
reports 2.27$\times$ inference geometric-mean speedup over 180+ models on one
A100 for TorchInductor, which generates Triton \citep{tillet2019triton};
KernelBench \citep{ouyang2025kernelbench} profiles 250 PyTorch workloads on an
L40S; TritonBench \citep{li2025tritonbench} evaluates 184 Triton operators on the
A100. Serving systems fix the kernel by preference order rather than by
measurement: vLLM \citep{kwon2023vllm}, Sarathi-Serve
\citep{agrawal2024sarathi} and DistServe \citep{zhong2024distserve} on A100,
Splitwise \citep{patel2024splitwise} on DGX-A100 and DGX-H100. TVM
\citep{chen2018tvm} and Ansor \citep{zheng2020ansor} instead search per target,
conceding that the right implementation is a property of the target. Outside ML
the question is mature: \citet{williams2009roofline} give the per-architecture
ceiling, and \citet{pennycook2016metric} and \citet{deakin2019portability}
define performance portability over a stated platform set and apply it across
twelve platforms.

\section{Study design}

\paragraph{Implementations and grid.} The prefill set is a naive reference, a
\code{torch.compile} path in three variants, the three PyTorch SDPA backends
(memory-efficient, flash, cuDNN), a hand-written Triton kernel, FlashAttention 2
and 3. The decode set is a naive KV-cache loop, an Inductor path,
SDPA, \code{flash\_attn\_with\_kvcache}, FlashInfer, and FlashAttention at query
length 1. A \emph{cell} keys regime, batch, query and KV head counts, head
dimension, sequence length, dtype, forward or forward-plus-backward, launch
mode, cache state and causality; the union grid is 159 prefill and 244 decode
cells, with per-device coverage in Table~\ref{tab:devices}.

\paragraph{Timing protocol.} Implementations are interleaved at the repetition
level inside one process, so a thermal or clock excursion lands on all of
them, not on whichever ran last. Latency is a CUDA-event median over inner
repetitions, then a median over process medians; confidence intervals come from
a 2000-draw cluster bootstrap over process medians, not over individual
repetitions, which share clocks and thermal point. The L2 cache
is flushed where the cell asks for a cold cache. Between-process CV is a median
0.011 on A10 (n=1871) and 0.026 on H100 (n=2474). Processes are pinned
to the NUMA node local to the GPU: two runs of the H100 decode grid at one
commit, differing only in pinning and restricted to the eager warm cells, give a
median between-process CV of 0.1116 unpinned against 0.0341 pinned over 966
(cell, implementation) triples spanning 192 cells on each side; the share of
those triples with a CV above 25\% falls from 30.7\% to 9.9\%. Both are timed by
the \code{block\_bench} wall-clock timer, not by CUDA events, so neither number
is comparable with the 0.026 above.

\paragraph{Correctness gate.} Every implementation is checked against a
reference in its equivalence class, and only rows that pass enter a ranking. Of 73{,}230 rows, 54{,}213 (74.0\%) ran and passed;
16{,}326 are \code{UNSUPPORTED}, 1{,}590 were refused in advance by an
out-of-memory predictor, and 950 raised a run-time error, all of them one
implementation on one device (Section~\ref{sec:mech}). The remaining 151 hit the
allocator despite passing the predictor, a defect in the predictor and not in
the implementation it was meant to protect. No row failed numerically. All
claims are within-device orderings compared across devices, never that device A
is faster than device B, and bandwidth denominators are measured on each host,
not taken from a datasheet.

\begin{table}[tb]
\centering
\small
\setlength{\tabcolsep}{3.2pt}
\begin{tabular}{lllrrr}
\toprule
GPU & cc & family & BW & reps & cells \\
 & & & (GB/s) & pre/dec & pre/dec \\
\midrule
A100     & 8.0  & Ampere    & 1732.8 & 5/5 & 144/192 \\
A10      & 8.6  & Ampere    &  483.3 & 5/5 & 144/191 \\
L40      & 8.9  & Ada       &  644.4 & 3/3 & 144/192 \\
L40S     & 8.9  & Ada       &  644.4 & 5/5 & 144/192 \\
H100     & 9.0  & Hopper    & 2979.7 & 5/5 & 159/244 \\
RTX 5090 & 12.0 & Blackwell & 1517.5 & 5/5 & 144/192 \\
\bottomrule
\end{tabular}
\caption{The six devices. BW is peak copy bandwidth measured in-process on that
host, never a vendor figure. \emph{reps} is the designed number of independent
process repeats per cell; where a run was cut short the realised median is
lower, at 2 for L40. \emph{cells} counts configurations
with at least one usable measurement.}
\label{tab:devices}
\end{table}

\section{Does the benchmark measure what it claims?}

A latency is evidence about a kernel only if that kernel is the one that ran, so
we record the launched kernel names for every timed row and compare them with
the implementation the harness asked for. Of the 54{,}213 passing rows, 36{,}054
(66.5\%) captured a readable trace. Grouping by the tuple that determines
dispatch, (implementation, GPU, head dimension, dtype, GQA ratio, mode, launch),
gives 651 dispatch classes, 625 of them (96.0\%) with at least one probed row.
The class is a backend-family unit and certifies only that: within a class the
template specialisation does vary with batch, and 156 of the flash-family
classes launch both \code{flash\_fwd\_splitkv\_combine} and \code{flash\_fwd},
while the family the row was asked for is the same in every one of them.
Class coverage is therefore not shape coverage, and the per-row figure above
is the one to read for that.
Among probed rows the requested kernel family is absent from the trace in 162
cases, a mismatch rate of 0.449\%, or 0.299\% over every audited row. The
unprobed 33.5\% count as unknown, never as agreement. Mismatches concentrate:
66 of the 162 are the SDPA memory-efficient backend on H100. The cells behind
the headline comparisons are probed at the dataset's own rate: of the 551
device-cells entering the restricted columns of Table~\ref{tab:flips}, 63.5\% of
passing rows carry a trace, and 51 of the 162 mismatches land on one of them.

The Inductor result needs its exact scope. Across 17{,}537 Inductor rows,
3{,}645 captured TorchInductor's \code{fuse\_attention} counter, and on all
3{,}645, on all six devices, that counter is 0. Inductor's own attention-fusion
pattern did not fire on any row whose counter we could read. The stronger claim,
that no Inductor row reached a flash kernel at all, does not hold: 45 probed
Inductor rows, all on H100, over 6 distinct cells, do carry a flash kernel in
their trace.

\begin{figure}[tb]
\centering
\includegraphics[width=\columnwidth]{fig6_winner_bands.pdf}
\caption{Fastest implementation, binned to workload bands: regime crossed with
batch (1, 4 to 8, 16 or more) and length (short $\le$1024, mid $\le$4096, long
above). Each column is one device, the colour names the implementation winning
most cells of that band there, and blank means the device ran none. A band is
filled only where at least half its cells have a winner separated by the
criterion of Section~\ref{sec:port}; the other 67 of the 99 non-empty
band-device cells are hatched, where the colour is a mode over cells that are
mostly ties, so a colour change across a hatched row is not evidence of a real
difference. All 2{,}082 (device, cell) pairs with two or more usable
implementations are binned, not only cells every device ran.}
\label{fig:winnermap}
\end{figure}

\section{Portability}
\label{sec:port}

\begin{figure*}[tb]
\centering
\includegraphics[width=\textwidth]{figA_separation.pdf}
\caption{How often a fastest backend is separated from the runner-up, by both a
10\% margin on the runner-up/fastest ratio and a cluster-bootstrap interval on
that ratio excluding 1. Counts are printed inside each bar; whiskers are 95\%
intervals over cells. \emph{Observation:} the criterion is met on 128 of
447 forward-prefill and 409 of 1203 decode configurations, and no device
clears half in both regimes. \emph{Interpretation:} exact winner counts are
sensitive to small timing differences over most of this grid. Configurations
that fail the criterion are not shown to be equal; they are configurations
where we could not establish a fastest backend.}
\label{fig:sep}
\end{figure*}

\begin{figure*}[tb]
\centering
\includegraphics[width=\textwidth]{figB_transfer.pdf}
\caption{Whether the selected backend transfers, and what it costs when it does
not. \emph{Left:} winner-flip rate over device pairs among Ampere, Ada and
Hopper, over each device's native backend set and over the set both devices of
a pair ran, with winners and separation re-derived after restricting.
\emph{Right:} the distribution of the target-device latency of the source
device's chosen backend divided by the target's own best.
\emph{Observation:} winners differ on 27 of 64 forward-prefill and 29
of 289 decode comparisons, yet the median transfer costs 1.002$\times$ and
1.000$\times$, with 95th percentiles of 1.61$\times$ and 1.32$\times$;
14\% of source choices do not run on the target at all and are counted here
rather than dropped. \emph{Interpretation:} a changed winner is usually not a
changed outcome. The cost of carrying a choice across these devices lives in
the tail and in availability, not at the median.}
\label{fig:transfer}
\end{figure*}


\paragraph{Most cells have no decisive winner.} A winner counts as \emph{separated} only if the runner-up's median latency is
at least 1.10 times the winner's and the cluster-bootstrap CI on that ratio
excludes 1. By that criterion Figure~\ref{fig:sep} gives the rate per device. Pooled, only 28.6\% of
forward-prefill and 34.0\% of decode configurations have a
separated winner (128 of 447 and 409 of 1203), and Figure~\ref{fig:winnermap} hatches every band where they
are a minority. A flip between implementations within a few percent of each
other is a coin toss reported as a result, so every rate below is computed
twice, over all shared cells and over cells separated on both devices.

\paragraph{The regime split.} Table~\ref{tab:flips} pools the ten device pairs
among A10, A100, L40, L40S and H100. Restriction pulls the two regimes apart. Forward prefill flips on 27 of 64 separated comparisons, interval [29, 53],
while decode flips on 29 of 289, interval [5, 15]. Restricting each device
pair to the backends both devices ran, and re-deriving winners and separation
over that restricted set rather than filtering afterwards, moves prefill to 33 of
72 and leaves decode at 29 of 289. Figure~\ref{fig:transfer} puts those rates beside what a flip costs: the
median transferred choice is 1.002$\times$ the target's own best in prefill
and 1.000$\times$ in decode, so most flips are between backends that are
practically interchangeable on the target. The restriction is what controls for
FlashAttention 3 existing only on H100 and FlashInfer being absent on the RTX
5090, and it does not rescue prefill.

\paragraph{RTX 5090 comparisons under the tested build.} The lower half of Table~\ref{tab:flips}
repeats the calculation for the five pairs including the RTX 5090. The decode rate stays high, at 63 of 85, while forward prefill is 16 of 30.
What makes the 5090 disagree is therefore specific to decode, and
Section~\ref{sec:mech} identifies it as availability, not architecture: in 62 of
those 64 flips the winner on the other device is FlashInfer, which produced no
timing on the 5090.

\begin{table}[tb]
\centering
\small
\setlength{\tabcolsep}{1.5pt}
\begin{tabular}{llrrrr}
\toprule
 & & \multicolumn{2}{c}{all shared cells} & \multicolumn{2}{c}{separated both} \\
\cmidrule(lr){3-4}\cmidrule(lr){5-6}
device set & regime & n & flip & n & flip [95\% CI] \\
\midrule
\multirow{2}{*}{Amp/Ada/Hop} & prefill & 720 & 66.0\% & 64 & 42.2\% [29, 53] \\
                             & decode  & 1916 & 52.1\% & 289 & 10.0\% [5, 15] \\
\midrule
\multirow{2}{*}{5090 vs rest} & prefill & 360 & 65.8\% & 30 & 53.3\% [30, 74] \\
                              & decode  & 959 & 78.4\% & 85 & 74.1\% [56, 90] \\
\bottomrule
\end{tabular}
\caption{Winner-flip rate, pooled over device pairs. \emph{Amp/Ada/Hop} pools
the ten pairs among A10, A100, L40, L40S and H100; \emph{5090 vs rest} the five
joining the RTX 5090 to each. \emph{n} counts device-pair comparisons on cells
both devices ran, \emph{flip} the share where the fastest implementation
differs. The right columns keep only cells separated on both devices; their
interval is a 95\% cluster bootstrap over cells, 2000 draws, since one cell
contributes to several pairs.}
\label{tab:flips}
\end{table}

\section{Mechanism: which binaries shipped}
\label{sec:mech}

\paragraph{The decode leader does not run on Blackwell.} FlashInfer
is the decode winner on 128 of 191 cells on A10, 133 of 192 on A100, 131 of 244
on H100 and 87 of 192 on L40, though only 49 of 192 on L40S, where the
SDPA path leads instead. On the RTX 5090 it wins nothing because it never
ran: of 960 attempted rows, 950 raised \code{RuntimeError: FlashInfer requires
GPUs with sm75 or higher} across 190 distinct cells, and the other 10 hit the
allocator. The text names sm75 on an sm120 part, so it is not a diagnosis, and
how much it would have won by there cannot be measured. What can be measured is
the fallback: the RTX 5090's decode winner is
\code{flash\_attn\_with\_kvcache} on 137 of 192 cells at a median 115.7 \us{}.

\paragraph{The wheels themselves.} Table~\ref{tab:prov} is a coverage map, not a
performance result. Its first two rows are the packaging asymmetry the prefill
flip rates sit beside:
FlashAttention 2 is the prefill winner on 43 of 144 A10 cells, 85 of 144 A100
cells and 72 of 144 RTX 5090 cells, and on H100 it wins none of the 159 while
FlashAttention 3 takes 81. That pattern is \emph{consistent with} the packaging
but not shown to be caused by it: we did not rebuild flash-attn for the missing
targets. FlashInfer is the other kind of gap: it JITs, so
its coverage belongs to the toolkit on the host, not to the wheel. We report its
sm\_120 failure and not a cause: every other library in the audit carries native
sm\_120 SASS built under the same CUDA 12.8, so a blanket claim that the toolkit
cannot target sm\_120 is contradicted by the audit's own rows.

\begin{table*}[tb]
\centering
\small
\setlength{\tabcolsep}{6pt}
\begin{tabular}{lllrrrr}
\toprule
backend & regime & GQA & support & OK & unsup. & fail \\
\midrule
D0-naive-kv & dec & both & \texttt{YYYYYY} & 5530 & 0 & 10 \\
D1-inductor & dec & expanded & \texttt{YYYYYY} & 194 & 5346 & 0 \\
D2-sdpa & dec & native & \texttt{YYYYYY} & 5535 & 0 & 5 \\
D3-fa-kvcache & dec & native & \texttt{YYYYYY} & 5515 & 0 & 25 \\
D4-flashinfer & dec & native & \texttt{YYYYY.} & 4565 & 0 & 975 \\
D6-fa-prefill-at-1 & dec & native & \texttt{YYYYYY} & 5515 & 0 & 25 \\
P0-naive & pre & both & \texttt{YYYYYY} & 3148 & 0 & 851 \\
P1-inductor & pre & both & \texttt{YYYYYY} & 3209 & 0 & 790 \\
P1-inductor-nofuse & pre & both & \texttt{YYYYYY} & 121 & 3873 & 5 \\
P1-inductor-where & pre & both & \texttt{YYYYYY} & 121 & 3873 & 5 \\
P2b-sdpa-mem-eff & pre & native & \texttt{YYYYYY} & 3969 & 30 & 0 \\
P2c-sdpa-flash & pre & native & \texttt{YYYYYY} & 3999 & 0 & 0 \\
P2d-sdpa-cudnn & pre & native & \texttt{YYYYYY} & 3999 & 0 & 0 \\
P3-triton & pre & both & \texttt{YYYYYY} & 3999 & 0 & 0 \\
P4-fa2 & pre & native & \texttt{YYYYYY} & 3999 & 0 & 0 \\
P4h-fa3 & pre & native & \texttt{..Y...} & 795 & 3204 & 0 \\
\bottomrule
\end{tabular}
\caption{Which comparisons are possible at all. \emph{support} is one character
per device in the order A10, A100, H100, L40, L40S, RTX 5090: \texttt{Y} produced at least one
usable row there, \texttt{.} attempted and produced none, \texttt{-} never attempted.
\emph{unsup.} counts configurations the backend declined in advance and \emph{fail}
runtime errors plus allocations refused or hit. FlashInfer reads \texttt{YYYYY.}: it ran
on five devices and produced no timing on the RTX 5090. FlashAttention 3 reads
\texttt{..Y...}, H100 only, by construction.}
\label{tab:avail}
\end{table*}



\section{Threats to validity}
\label{sec:threats}

\paragraph{Decode spans two harness revisions, and one host per device.} The
A10 and A100 decode rows carry harness commit \code{468b9a55} and every other
decode row \code{91959d34}, recorded days later. Diffing the two revisions
shows the measurement path is unchanged: \code{bench.py}, \code{impls.py},
\code{check.py} and the vendored Triton kernel are byte-identical, the image
pins the same library versions, and the \code{run.py} differences decide which
cells are attempted rather than how a cell is timed. The revision label
therefore does not mark a change in what was measured, and we do not treat it
as one. Driver and host are the real covariates, and they cannot be separated
from device in this design: each device was measured on the host that had it.
Neither label predicts disagreement. The L40 rows share driver 595.71.05 with
A10 and A100 yet carry the later commit, and A10 against L40 flips on 0 of 36
separated comparisons while H100 against L40, one commit and two drivers apart,
flips on 9 of 33. What predicts the decode flip rate is which architectures are
being compared, which is the decomposition reported in Section~\ref{sec:port}.
A single-host replication of the whole grid would still be worth having, and we
did not have six identical hosts. Prefill is clean in that sense: all six devices ran prefill at one commit.
A100 prefill was however re-measured on a rented host, driver 580.126.16,
after the cluster run reached only one usable process repeat, so that device's
two regimes come from two hosts. No comparison crosses the boundary, since
every ranking is formed within one device and one regime, and the two hosts
measured copy bandwidth within 0.05\% of each other.

\paragraph{Thin repeats, and an expensive restriction.} L40 has 3 designed
process repeats against 5 on every other device, and a realised median of 2,
though its cell coverage is complete. Across the dataset 1 of 11{,}858 usable
triples has a single process repeat, and for it the bootstrap has nothing to
resample, so its CI width of 0 is an artefact. To check that thinner
sampling does not manufacture separated winners, we recount separation on
two complete devices after subsampling them to three repeats: L40S moves
from 31 to 34 separated decode cells and A100 from 115 to 110, so the
criterion is dominated by the 1.10 margin rather than by bootstrap width. Requiring a separated winner on
both sides then discards about two thirds of the cells as ties, which is the
right call for a ranking claim but leaves the headline comparisons resting on
tens to low hundreds of paired cells.

\paragraph{Unprobed dispatch cells.} 33.5\% of passing rows captured no kernel
trace and 26 of 651 dispatch classes have no probed row at all, so the 0.449\%
mismatch rate is a rate among probed rows. On the unprobed rows the answer is
unknown, not zero.

\paragraph{Two floors under the measurement.} On L40 the median ratio of
effective KV traffic to that host's measured peak copy bandwidth is 1.038, so
the peak probe under-reads what the decode kernels sustain there. No L40 number
should be read as a fraction of achievable bandwidth: the effective GB/s
stands, the ratio does not. On H100 the fastest usable cell is 6.34 \us{} and 153 of its 2{,}474 triples sit
within 5$\times$ of the 5.088 \us{} CUDA-event pair overhead. That comparison
crosses timers and we do not draw a floor from it: 152 of those 153 triples are
\code{do\_bench\_cudagraph} cells, while 5.088 \us{} is a \code{block\_bench}
event-pair cost. What the two together do establish is that at the small-batch,
short-context, CUDA-graph end of the decode grid the winner margins are of the
same order as the timing machinery, and we have not measured that end well
enough to rank it. Those cells are reported and are not used to support any
portability claim.

\paragraph{What was not measured.} Every cloud host we rented blocked
Nsight Compute and Nsight Systems, so no occupancy, cache-hit-rate or
counter-derived roofline claim appears here, and every mechanism in
Section~\ref{sec:mech} rests on a static binary audit and on run-time status
codes. We measured exactly one torch, flash-attn and flashinfer combination, so a
different wheel is a different experiment rather than a replication. What is
timed is a single attention call on six devices from one vendor.

\section{Conclusion}

On these six devices, which attention implementation is fastest is not a
property of the implementation, and a benchmark naming none of the device, the
regime and the distribution has under-specified its result. Two practical
consequences follow. A serving system that fixes its attention backend by
preference order should re-measure that order on each deployment target, since
the orderings that transfer here are the decode ones, and only within the
architectures for which the installed libraries ship native code. And a
reported winner should carry the margin separating it from the runner-up,
because on roughly two thirds of the configurations measured here there is no
separation to report.

\appendix
\section{Binary provenance}
\label{app:prov}
The audit records which architectures the installed wheels carry native code
for. It establishes artifact coverage, not a performance cause: no rebuild for
a missing target was performed, and CUDA permits a cubin built for a lower minor
compute capability to run on a higher one within a major version.

\begin{table}[tb]
\centering
\small
\setlength{\tabcolsep}{3pt}
% Seven columns do not fit a column at \small; scale rather than shrink the
% font, so this table keeps the same look as the other two.
\resizebox{\columnwidth}{!}{%
\begin{tabular}{lcccccc}
\toprule
library & sm\_80 & sm\_86 & sm\_89 & sm\_90 & sm\_120 & PTX \\
\midrule
flash-attn 2.8.3   & SASS & \emph{80} & \emph{80} & SASS & SASS & none \\
libtorch\_cuda     & SASS & SASS & SASS & SASS & SASS & none \\
cuDNN engines      & SASS & SASS & \emph{86} & SASS & SASS & sm\_120 \\
cuBLAS             & SASS & SASS & \emph{86} & SASS & SASS & sm\_120 \\
flashinfer 0.6.18  & JIT & JIT & JIT & JIT & none & n/a \\
flash-attn 3       & n/f & n/f & n/f & n/f & n/f & n/f \\
\bottomrule
\end{tabular}}
\caption{Binary provenance of the installed wheels, from \code{cuobjdump}, over
the compute capabilities of the six devices. SASS means native machine code is
present; \emph{PTX} is the embedded PTX a driver could compile where there is no
cubin. An italic number means none was shipped and the target runs that older
architecture's cubin under CUDA minor-version binary compatibility. JIT means
the library ships no cubins and builds at run time, \emph{none} that this build
produced no kernel there under CUDA 12.8, and \emph{n/f} that the audit found no
such library, which for FlashAttention 3, a path that ran on H100, is a gap in
the audit and not in the install.}
\label{tab:prov}
\end{table}

\bibliography{custom}

\end{document}
